\documentclass[10pt]{extarticle}
\usepackage[a4paper,margin=1.5cm]{geometry}
\usepackage[T1]{fontenc}
\usepackage{amsmath,amssymb}
\usepackage{graphicx}
\usepackage[export]{adjustbox}
\usepackage{multicol}
\usepackage{enumitem}
\usepackage{booktabs,array,tabularx}
\setlength{\columnsep}{0.7cm}
\setlength{\parindent}{0pt}
\setlength{\parskip}{3pt}
\pagestyle{empty}
\makeatletter
\newcommand{\qfigdef}[4]{\@namedef{qf@file@#1}{#2}\@namedef{qf@w@#1}{#3}\@namedef{qf@wide@#1}{#4}}
\newcommand{\qf@check}[1]{\@ifundefined{qf@file@#1}{\PackageError{score2}{Unknown figure id #1}{}}{}}
\newcommand{\qf@img}[4][]{\includegraphics[width=#2\dimexpr\@nameuse{qf@w@#4}pt\relax,max width=#3,max height=0.45\textheight,#1]{\@nameuse{qf@file@#4}}}
\newif\ifqf@full
\newif\ifqf@boxed
\newsavebox{\qf@box}
\newcommand{\qf@open}{\global\qf@fulltrue\par\global\qf@boxedtrue
  \setbox\qf@box\vbox\bgroup\hsize\textwidth\linewidth\textwidth}
\newcommand{\qf@emit}[1]{\par\egroup\global\qf@boxedfalse
  \ifdim\dimexpr\ht\qf@box+\dp\qf@box\relax<#1\textheight\noindent\usebox{\qf@box}\par\else\unvbox\qf@box\fi}
\newcommand{\qfullsplit}{\ifqf@boxed\qf@emit{0.85}\fi}
\newcommand{\qf@close}{\ifqf@boxed\qf@emit{0.45}\fi\global\qf@fullfalse}
\newcommand{\qfullbegin}{\end{multicols}\qf@open}
\newcommand{\qfullend}{\qf@close\begin{multicols}{2}}
\newcommand{\qfullstart}{\qf@open}
\newcommand{\qfulldone}{\qf@close}
\newcommand{\qfig}[2][1]{\qf@check{#2}%
  \ifqf@full\qf@blk{#1}{#2}\else
  \ifnum\@nameuse{qf@wide@#2}=1 \expandafter\@firstoftwo\else\expandafter\@secondoftwo\fi
  {\end{multicols}\begin{center}\qf@img{#1}{\textwidth}{#2}\end{center}\begin{multicols}{2}}%
  {\qf@blk{#1}{#2}}\fi}
\newcommand{\qf@blk}[2]{\par\smallskip{\centering\qf@img{#1}{\linewidth}{#2}\par}\smallskip}
\newcommand{\qfigi}[2][1]{\qf@check{#2}\qf@img[valign=c]{#1}{\linewidth}{#2}}
\makeatother
\newcommand{\papersec}[1]{\par\bigskip{\centering\Large\bfseries PAPER-#1\par}\medskip}
\newcommand{\partsec}[2]{\par\medskip\noindent{\large\bfseries PART-#1: #2}\par\smallskip}
\newcommand{\sect}[1]{\par\smallskip\noindent\textbf{#1}\par}
\newcommand{\qn}[1]{\par\smallskip\noindent\textbf{#1.}~}
\newcommand{\paraq}[2]{\par\smallskip\noindent\textbf{Paragraph for Questions #1 and #2}\par}
\newenvironment{qoptions}{\begin{itemize}[leftmargin=1.8em,itemsep=1pt,topsep=2pt,label={}]}{\end{itemize}}
\newcommand{\opt}[2]{\item[(#1)] #2}
% figure registry (generated by score2_tools.py finalize)
\qfigdef{p1_phy_s1_q05_1}{Score-I - Test 07 (04 Feb 2022)-img/p1_phy_s1_q05_1.png}{141.3}{0}
\qfigdef{p1_phy_s2_q02_1}{Score-I - Test 07 (04 Feb 2022)-img/p1_phy_s2_q02_1.png}{207.2}{0}
\qfigdef{p1_phy_s2_q03_1}{Score-I - Test 07 (04 Feb 2022)-img/p1_phy_s2_q03_1.png}{38.4}{0}
\qfigdef{p1_phy_s2_q05_1}{Score-I - Test 07 (04 Feb 2022)-img/p1_phy_s2_q05_1.png}{69.3}{0}
\qfigdef{p1_phy_s3_q01_1}{Score-I - Test 07 (04 Feb 2022)-img/p1_phy_s3_q01_1.png}{104.9}{0}
\qfigdef{p1_phy_s3_q03_1}{Score-I - Test 07 (04 Feb 2022)-img/p1_phy_s3_q03_1.png}{122.7}{0}
\qfigdef{p1_phy_s3_q04_1}{Score-I - Test 07 (04 Feb 2022)-img/p1_phy_s3_q04_1.png}{98.5}{0}
\qfigdef{p1_phy_s3_q05_1}{Score-I - Test 07 (04 Feb 2022)-img/p1_phy_s3_q05_1.png}{81.1}{0}
\qfigdef{p1_phy_s3_q06_1}{Score-I - Test 07 (04 Feb 2022)-img/p1_phy_s3_q06_1.png}{146.9}{0}
\qfigdef{p1_chem_s1_q03_1}{Score-I - Test 07 (04 Feb 2022)-img/p1_chem_s1_q03_1.png}{72.7}{0}
\qfigdef{p1_chem_s1_q04_A}{Score-I - Test 07 (04 Feb 2022)-img/p1_chem_s1_q04_A.png}{74.0}{0}
\qfigdef{p1_chem_s1_q04_B}{Score-I - Test 07 (04 Feb 2022)-img/p1_chem_s1_q04_B.png}{70.5}{0}
\qfigdef{p1_chem_s1_q04_C}{Score-I - Test 07 (04 Feb 2022)-img/p1_chem_s1_q04_C.png}{93.1}{0}
\qfigdef{p1_chem_s1_q04_D}{Score-I - Test 07 (04 Feb 2022)-img/p1_chem_s1_q04_D.png}{146.7}{0}
\qfigdef{p1_chem_s2_q02_1}{Score-I - Test 07 (04 Feb 2022)-img/p1_chem_s2_q02_1.png}{52.5}{0}
\qfigdef{p1_chem_s2_q03_1}{Score-I - Test 07 (04 Feb 2022)-img/p1_chem_s2_q03_1.png}{114.4}{0}
\qfigdef{p1_chem_s2_q04_1}{Score-I - Test 07 (04 Feb 2022)-img/p1_chem_s2_q04_1.png}{150.6}{0}
\qfigdef{p1_chem_s2_q05_i}{Score-I - Test 07 (04 Feb 2022)-img/p1_chem_s2_q05_i.png}{47.7}{0}
\qfigdef{p1_chem_s2_q05_ii}{Score-I - Test 07 (04 Feb 2022)-img/p1_chem_s2_q05_ii.png}{53.7}{0}
\qfigdef{p1_chem_s2_q05_iii}{Score-I - Test 07 (04 Feb 2022)-img/p1_chem_s2_q05_iii.png}{73.7}{0}
\qfigdef{p1_chem_s2_q05_iv}{Score-I - Test 07 (04 Feb 2022)-img/p1_chem_s2_q05_iv.png}{75.2}{0}
\qfigdef{p1_chem_s2_q05_v}{Score-I - Test 07 (04 Feb 2022)-img/p1_chem_s2_q05_v.png}{57.8}{0}
\qfigdef{p1_chem_s2_q05_vi}{Score-I - Test 07 (04 Feb 2022)-img/p1_chem_s2_q05_vi.png}{54.6}{0}
\qfigdef{p1_chem_s2_q05_vii}{Score-I - Test 07 (04 Feb 2022)-img/p1_chem_s2_q05_vii.png}{54.8}{0}
\qfigdef{p1_chem_s2_q05_viii}{Score-I - Test 07 (04 Feb 2022)-img/p1_chem_s2_q05_viii.png}{25.0}{0}
\qfigdef{p1_chem_s2_q05_ix}{Score-I - Test 07 (04 Feb 2022)-img/p1_chem_s2_q05_ix.png}{56.8}{0}
\qfigdef{p1_chem_s2_q05_x}{Score-I - Test 07 (04 Feb 2022)-img/p1_chem_s2_q05_x.png}{53.7}{0}
\qfigdef{p1_chem_s2_q06_1}{Score-I - Test 07 (04 Feb 2022)-img/p1_chem_s2_q06_1.png}{99.6}{0}
\qfigdef{p1_chem_s3_q01_g}{Score-I - Test 07 (04 Feb 2022)-img/p1_chem_s3_q01_g.png}{36.3}{0}
\qfigdef{p1_chem_s3_q02_a}{Score-I - Test 07 (04 Feb 2022)-img/p1_chem_s3_q02_a.png}{32.5}{0}
\qfigdef{p1_chem_s3_q02_b}{Score-I - Test 07 (04 Feb 2022)-img/p1_chem_s3_q02_b.png}{26.3}{0}
\qfigdef{p1_chem_s3_q02_d}{Score-I - Test 07 (04 Feb 2022)-img/p1_chem_s3_q02_d.png}{26.5}{0}
\qfigdef{p1_chem_s3_q03_1}{Score-I - Test 07 (04 Feb 2022)-img/p1_chem_s3_q03_1.png}{262.0}{1}
\qfigdef{p1_math_s2_q03_1}{Score-I - Test 07 (04 Feb 2022)-img/p1_math_s2_q03_1.png}{72.0}{0}

\begin{document}
\begin{center}{\Large\bfseries JEE (Main+Advanced) Score-I --- Test 07 (04 Feb 2022)}\end{center}
\begin{multicols}{2}
\raggedcolumns
\partsec{1}{PHYSICS}
\sect{Section-I}
\qn{1} A travelling wave pulse on a string of linear mass density $\dfrac{3\,\mathrm g}{\mathrm{cm}}$ is given by $y=2.00e^{-\left(4x^2+324t^2+72xt+3\right)}$ where x and y are in meter and t in second then choose the correct options.
\begin{qoptions}
\opt{A}{Wave pulse is travelling along positive x-axis with speed 9 m/s.}
\opt{B}{Wave pulse is travelling along negative x-axis with speed 9 m/s.}
\opt{C}{Tension in the string is 24.3 N}
\opt{D}{Amplitude of wave pulse is 2.00 m}
\end{qoptions}
\qn{2} A wave disturbance in a medium is described by $y(x,t)=0.08\cos\left(100\pi t+\dfrac{\pi}{2}\right)\cos(20\pi x)$ where x and y are in meter and t in second.
\begin{qoptions}
\opt{A}{A node occurs at x = 0.125 m}
\opt{B}{An antinode occurs at x = 0.3 m}
\opt{C}{Speed of component wave is 5 m/s}
\opt{D}{The wavelength is 0.2 m}
\end{qoptions}
\qn{3} Two observers O$_1$ and O$_2$ are at distance r and 10r from a point source of sound wave
\begin{qoptions}
\opt{A}{If the intensity of sound received by O$_1$ is 40dB, the intensity received by O$_2$ will be 20dB}
\opt{B}{If the intensity received by O$_1$ is 100 dB, the intensity received by O$_2$ will 80 dB}
\opt{C}{If the intensity received by O$_1$ is $10^{-6}\,\dfrac{\mathrm w}{\mathrm m^2}$, the intensity received by O$_2$ will be $10^{-8}\,\dfrac{\mathrm w}{\mathrm m^2}$}
\opt{D}{If the intensity received by O$_1$ is $10^{-10}\,\dfrac{\mathrm w}{\mathrm m^2}$, the intensity received by O$_2$ will be $10^{-12}\,\dfrac{\mathrm w}{\mathrm m^2}$}
\end{qoptions}
\qn{4} A source emitting a sinusoidal sound wave of frequency 360 Hz travels towards a wall at a speed 30 m/s. What is the beat frequency perceived by an observer moving at speed 10 m/s, speed of sound in air is 330 m/s ?
\begin{qoptions}
\opt{A}{64Hz if observer is going away from the source and wall.}
\opt{B}{68 Hz if observer is going towards the source and the wall}
\opt{C}{60 Hz if observer is in between the wall and the source and moving away from the wall.}
\opt{D}{24 Hz if observer is in between the wall and the source, and moving towards the wall.}
\end{qoptions}
\qn{5} In YDSE a screen is at a distance 80 cm from two narrow slits s$_1$ and s$_2$ which are 2 mm apart. Slit s$_1$ is covered by a transparent sheet of thickness 2.5 $\mu$m and slit s$_2$ is covered by another sheet of thickness 1.25 $\mu$m as shown.
\qfig{p1_phy_s1_q05_1}
Both sheets are made of same material of refractive index 1.4. Water is also filled between slits and screen. Then [Wavelength of light in air = 5000\AA, Refractive index of water = 4/3]
\begin{qoptions}
\opt{A}{Central maxima is above the point C}
\opt{B}{Distance of central maxima is 0.015mm from C.}
\opt{C}{The resultant intensity at C is 75\% of maximum intensity}
\opt{D}{Fringe width is 0.15 mm.}
\end{qoptions}
\qn{6} A particle performing SHM is found at its equilibrium position at t = 1 sec and it's velocity (in positive direction) is 0.25 m/s at t = 2 sec. If the time period of oscillation is 6 sec. Then--
\begin{qoptions}
\opt{A}{Amplitude of oscillation is $\dfrac{3}{2\pi}$m}
\opt{B}{Initial phase is $\left(-\dfrac{\pi}{3}\right)$}
\opt{C}{Velocity of particle at 6 sec is 0.25 m/s}
\opt{D}{Acceleration of particle at 6 sec is zero.}
\end{qoptions}
\sect{Section-II}
\qn{1} The position of a particle of mass 1 kg with respect to origin varies according to relation $x=3\sin(100t)+8\cos^2 50t$, where x is in meter and t in sec. Find out total energy of particle in (kJ) \{Assume PE at mean position is zero\}
\qn{2} Stationary wave is setup in a uniform string clamped at both ends. Length of string is 0.6 m. Snapshots of string are taken at two instant one at t = 0 sec and another at t = 0.4 sec. These two snapshots are shown below. Velocity of point P (which is also the midpoint of string) is in upward direction (take upward direction positive) at t = 0 sec. At the instant snapshots are taken particles are half of their respective maximum displacement from mean position. During this time interval particles have crossed their mean position only once. Find out velocity (in m/s) of travelling wave in string.
\qfig{p1_phy_s2_q02_1}
\qn{3} A uniform inextensible rope of length 20m and a light inextensible thread are connected at their ends A and B to move a loop. The thread passes over a fixed ideal pulley and the rope passes under another fixed ideal pulley. The junction B is pulled slightly downwards and then released. Find period of oscillations (in sec). (Use g = 10 m/s$^2$ and $\pi$ = 3.14)
\qfig{p1_phy_s2_q03_1}
\qn{4} A string fixed at both ends with tension 100N is vibrating in lowest mode of vibration for which a point at quarter of its length from one end is a point of maximum displacement. The frequency of vibration in this mode is 30 Hz. The frequency (in Hz) emitted when it vibrates in the next mode with new tension 25N such that this point is again a point of maximum displacement is
\qn{5} In a standard Young's double-slit experiment, a thin layer of air of thickness t$_2$ is enclosed in a hollow squarish transparent container of wall thickness t$_1$ as shown in figure.
\qfig{p1_phy_s2_q05_1}
When it is placed in front of one slit, causes the shifting of central maxima by 4 fringe width, and when the entire setup is immersed in a liquid then it is observed that the central maxima is restored to the original position. Find the side length of container in $\mu$m.

[Wavelength of light in air = 500 nm, Refractive index of container = 1.5, Refractive index of liquid = 1.4]
\qn{6} Calculate the sound intensity level in decibel for a sound wave travelling in air at 0$^\circ$C and having a pressure amplitude of 0.5 Pa. Assume density of air = 1.25 kg/m$^3$, velocity of sound = 330 m/s, $\log_{10}(3.3)=0.52$
\sect{Section-III}
\qn{1} A thick rope of length 4L and linear mass density 16$\mu$ is joined at B to a thin rope of length 2L and linear mass density 4$\mu$. The system is horizontally stretched by two vertical wall A and C. Assuming B to be a node find minimum number of loops in thick rope
\qfig{p1_phy_s3_q01_1}
\qn{2} A 5.0 m long vertical pipe is filled completely with a liquid. A small hole is drilled at the base of pipe through which liquid starts leaking out. This pipe resonates with a tuning fork. The first time resonance occur when height of water column is 4.4 m and 2.8 m respectively. If radius of pipe is found to be $\left(\dfrac{1}{x}\right)$ m, find x
\qn{3} Figure shows arrangement of a sound source at (0, 1) mm and two detectors D$_1$ and D$_2$ placed at $\left(\mathrm D,\dfrac{d}{2}\right)$ and $\left(\mathrm D,-\dfrac{d}{2}\right)$ respectively. The source emitted a sound pulse at a certain time. Assuming velocity of sound in surrounding medium is 300 m/s. Time gap between the recordings made by D$_1$ and D$_2$ is $\dfrac1n\mu$s, find $\left(\dfrac{n}{30}\right)$ [take D = 1 m, d = 2 mm]
\qfig{p1_phy_s3_q03_1}
\qn{4} A torsional pendulum consists of a uniform disc of mass m and radius R attached to a thin rod of torsional constant C. Two particles of mass m strike with the disc with speed v each and stick to it as shown in figure.
\qfig{p1_phy_s3_q04_1}
Find the angular amplitude (radian) of oscillation.

[Take m = 2kg, v = 20 m/s, R = 20 cm, C = 80Nm/rad]
\qn{5} A thin film of refractive index $\sqrt3$ is coated on a diamond plate of refractive index 2.42 as shown. A parallel beam of white light falls from air on thin film at angle of incidence $60^\circ$. Find the minimum thickness (t) in $\mu$m if reflected light is most intense for $\lambda$ = 600 nm. Fill 10t.
\qfig{p1_phy_s3_q05_1}
\qn{6} A glass prism $\left(\mu_g=\dfrac32\right)$ of refracting angle $2^\circ$ is placed between alcohol ($\mu$ = 1.36) \& water ($\mu$ = 4/3) as shown in figure.
\qfig{p1_phy_s3_q06_1}
A monochromatic light of wavelength (in air) 500 nm falls normally on prism from alcohol side and form the interference pattern on a screen placed at 100 cm from prism. Calculate the fringe width in $\mu$m. (Take $\pi\simeq3$). Fill $\dfrac{\mathrm w}{5}$.
\raggedcolumns
\partsec{2}{CHEMISTRY}
\sect{Section-I}
\qn{1} Highly pure dilute solution of sodium in liquid ammonia.
\begin{qoptions}
\opt{A}{On evaporation yield metals}
\opt{B}{Exhibits electrical conductivity}
\opt{C}{Produces sodium amide and hydrogen gas instantly}
\opt{D}{Has much greater volume than the sum of volume of metal and solvent}
\end{qoptions}
\qn{2} Identify the correct order for solubility of 2$^\text{nd}$ group.
\begin{qoptions}
\opt{A}{Solubility of oxalates increases from Ca to Ba}
\opt{B}{Solubility of fluorides decreases from Ca to Ba}
\opt{C}{Solubility of hydroxides increases from Ca to Ba}
\opt{D}{Solubility of carbonates decrease from Ca to Ba}
\end{qoptions}
\qn{3} Incorrect statement among the following for the given compound `X' is/are
\qfig{p1_chem_s1_q03_1}
\begin{qoptions}
\opt{A}{Cyclo is used as primary prefix in the IUPAC nomenclature of `X'}
\opt{B}{Correct IUPAC nomenclature is 3-Aminobutane-1, 4-dioic anhydride}
\opt{C}{Correct IUPAC nomenclature is 3-Amino cyclobutanic-1, 4-dioic anhydride}
\opt{D}{Acid anhydride is the principal function group in `X' compound}
\end{qoptions}
\qn{4} Correctly matched among the following is/are

\hspace{1em}Column-I \hfill Column-II \hspace{1em}
\begin{qoptions}
\opt{A}{\qfigi{p1_chem_s1_q04_A} $-$ Heat of hydrogenation}
\opt{B}{\qfigi{p1_chem_s1_q04_B} $-$ Heat of combustion}
\opt{C}{\qfigi{p1_chem_s1_q04_C} $-$ Dipole moment}
\opt{D}{\qfigi{p1_chem_s1_q04_D} $-$ Stability}
\end{qoptions}
\qn{5} Which of the following is/are correct?
\begin{qoptions}
\opt{A}{Photocurrent depends on the intensity of incident photons}
\opt{B}{For more intense incident radiation, value of stopping potential is greater}
\opt{C}{Stopping potential for given metal depends only on the frequency of incident radiation}
\opt{D}{Slope of stopping potential v/s frequency graph is constant for all metals and intercept depends on the frequency of incident radiation}
\end{qoptions}
\qn{6} A reaction $\mathrm{A_2+B_2\rightleftharpoons2AB}$ follows the path :

$\mathrm{A_2}\ \overset{\text{fast}}{\rightleftharpoons}\ \mathrm{A+A}$

$\mathrm{A+B_2}\ \xrightarrow{\text{slow}}\ \mathrm{AB+B}$

$\mathrm{A+B}\ \overset{\text{fast}}{\rightleftharpoons}\ \mathrm{AB}$
\begin{qoptions}
\opt{A}{overall order of the reaction is 1.5}
\opt{B}{order w.r.t. A$_2$ is 0.5}
\opt{C}{order w.r.t. B$_2$ is 0.5}
\opt{D}{Order of a reaction may change by changing conditions like concentration, temperature or catalyst.}
\end{qoptions}
\sect{Section-II}
\qn{1} The uncertainty in the measurement of position of an electron in 2$^\text{nd}$ excited state of He$^+$ is $\dfrac{9}{16}$\AA. What will be the minimum uncertainty in the measurement of De-Broglie wave length of this electron in angstorms(\AA) ? [take first Bohr radius = 0.5 \AA] [$\pi$ = 3.14]
\qn{2} For parallel reaction of reactant X to products Y \& Z given :
\qfig{p1_chem_s2_q02_1}

The ratio $\dfrac{\mathrm k_1}{\mathrm k_2}=\dfrac23$ [k = rate constant]

When all of X is completely consumed then find mole \% of Z obtained?
\qn{3}
\qfig{p1_chem_s2_q03_1}
Number of nonaxial d-orbital in hybridisation of central metal in gas A = x

Number of Cr -- O bonds (sigma) in B = y

Molar mass of gas C = z

Find $\dfrac{z}{x+y}$ =?
\qn{4}
\qfig{p1_chem_s2_q04_1}
If in the above reaction scheme difference of oxidation number of metal between (B) \& (C) is equal to x and molecular mass of gas G is y. Find $\dfrac{y}{x}$.
\qn{5} Compounds having pKa value less than that of carbonic acid (i.e., H$_2$CO$_3$ 6.4 approx.) is X and greater than that of carbonic acid is Y. Then find out $\dfrac{\mathrm X-\mathrm Y}{3}$ =?

(i) \qfigi{p1_chem_s2_q05_i}

(ii) \qfigi{p1_chem_s2_q05_ii}

(iii) \qfigi{p1_chem_s2_q05_iii}

(iv) \qfigi{p1_chem_s2_q05_iv}

(v) \qfigi{p1_chem_s2_q05_v}

(vi) \qfigi{p1_chem_s2_q05_vi}

(vii) \qfigi{p1_chem_s2_q05_vii}

(viii) \qfigi{p1_chem_s2_q05_viii}

(ix) \qfigi{p1_chem_s2_q05_ix}

(x) \qfigi{p1_chem_s2_q05_x}
\qn{6} Degree of unsaturation in the given compound is = x.
\qfig{p1_chem_s2_q06_1}
Number of different types of functional group present = y. Then find out $\dfrac{2x+y}{y}$ =?
\sect{Section-III}
\qn{1} How many groups will show resonance with phenyl ring when attached with it?

(a) $-\mathrm{CH_2-CH=CH_2}$ \quad (b) $-\mathrm{CH_2-O-\overset{\mathrm O}{\overset{\|}{C}}-NH_2}$

(c) $-\overset{\oplus}{\mathrm N}\mathrm H_3$ \quad (d) $-\mathrm{N=O}$

(e) $-\mathrm{Cl}$ \quad (f) $-\overset{\oplus}{\mathrm C}\mathrm H_2$

(g) \qfigi{p1_chem_s3_q01_g} \quad (h) $-\mathrm{BH_2}$
\qn{2} How many compounds given below have C = C bond length shorter than $\mathrm{H_2C\underset{\uparrow}{=}CH-CH_3}$.

(a) \qfigi{p1_chem_s3_q02_a} \quad (b) \qfigi{p1_chem_s3_q02_b}

(c) $\mathrm{O_2N-CH\underset{\uparrow}{=}CH-OEt}$ \quad (d) \qfigi{p1_chem_s3_q02_d}

(e) $\mathrm{CH_3-CH_2-CH\underset{\uparrow}{=}CH_2}$

(f) $\mathrm{H_3C-CH\underset{\uparrow}{=}CH-CH_3}$

(g) $\mathrm{(CH_3)_2C\underset{\uparrow}{=}C(CH_3)_2}$
\qfullbegin
\qn{3}
\qfig{p1_chem_s3_q03_1}
\qfullsplit
If molecular mass of gas C, gas D and gas G are x, y \& z respectively. Find x -- y -- 2z.
\qfullend
\qn{4} Aqueous NaCl solution on electrolysis forms three products X, Y \& Z. X \& Y are gases. X can react with Z but Y cannot. X reacts with Z and can form two different oxo anion products P and Q under different temperature and varying conc. conditions. The sum of oxidation state of Cl in P and Q is :
\qn{5} The electron in Be$^{+3}$ ion makes a transition from n = 16 to n = 4. This photon is used to excite H-atom in ground state. The De-Broglie wave length of this excited electron in terms of `$\pi$\AA' is: (use first Bohr orbit of H = 0.5 \AA)
\qn{6} A reaction follows temperature - rate constant relation as

$\ell\mathrm n\,\mathrm k=30-\dfrac{5000}{\mathrm T}+3\,\ell\mathrm n(3\mathrm T)$

If activation energy = 11.8 kcal/mol and R = 2 cal k$^{-1}$ mol$^{-1}$ then find the value of $\dfrac{\mathrm T}{100}$ where T = reaction temperature [in kelvin].
\raggedcolumns
\partsec{3}{MATHEMATICS}
\sect{Section-I}
\qn{1} If $e^xf(x)=2(f(x))^2-3x$ and $g(x)=\left(\dfrac{e^x\left(2e^{2x}-8e^xf(x)+f(x)\right)+3}{\left(e^x-4f(x)\right)\left(e^{2x}-f(x)\right)^2}\right)$ where $f(x)>0\ \forall\ x\in\mathbb{R}$, then which of the following is/are true ? (Where c is constant of integration)
\begin{qoptions}
\opt{A}{$\displaystyle\int g(x)\,dx=\dfrac{-1}{e^x-f(x)}+c$}
\opt{B}{$\displaystyle\int g(x)\,dx=\dfrac{-1}{e^{2x}-f(x)}+c$}
\opt{C}{$\displaystyle\int_0^1 g(x)\,dx=2\left(\dfrac{2}{e+\sqrt{e^2+24-4e^2}}+1\right)$}
\opt{D}{$\displaystyle\int_0^1 g(x)\,dx=2\left(\dfrac{2}{e+\sqrt{e^2+24-4e^2}}+2\right)$}
\end{qoptions}
\qn{2} Let $\mathrm I_1=\displaystyle\int_0^{\ln10^4}\left(e^{2x}-e^x\left[e^x\right]\right)dx$ and $\mathrm I_2=\displaystyle\int_0^{9\times10^6}\left(1-\dfrac{\left[\sqrt x\right]}{\sqrt x}\right)dx$. Then which of the following is/are correct.

where [$\cdot$] denotes greatest integer function.
\begin{qoptions}
\opt{A}{$\mathrm I_1>\mathrm I_2$}
\opt{B}{$\dfrac{\mathrm I_1}{\mathrm I_2}=\dfrac35$}
\opt{C}{$\left[\dfrac{\mathrm I_1-\mathrm I_2}{1000}\right]=1$}
\opt{D}{$\dfrac{\mathrm I_1}{\mathrm I_2}=\dfrac53$}
\end{qoptions}
\qn{3} Let S$_k$ be the area bounded by the curve $y=x^2(1-x)^k$ and the lines $x=0$, $y=0$ and $x=1$. Then which of the following is/are correct.
\begin{qoptions}
\opt{A}{$\displaystyle\sum_{k=1}^{100}\mathrm S_k<\dfrac16$}
\opt{B}{If $\displaystyle\lim_{n\to\infty}\sum_{k=1}^{n}\mathrm S_k=\dfrac{p}{q}$, then $p+q=7$ where p and q are coprime numbers}
\opt{C}{$\mathrm S_{99}=\dfrac{1}{100\times101\times102}$}
\opt{D}{$\displaystyle\sum_{k=1}^{999}\mathrm S_k>\dfrac16$}
\end{qoptions}
\qn{4} Let $\displaystyle\int\sqrt{\tan x}\,e^{\sqrt{\cos x}}\sqrt{\sin x}\,dx=f(x)+C$ where $f(0)=-2e$ (where C is constant of integration), then -
\begin{qoptions}
\opt{A}{$f\left(\dfrac\pi4\right)=-2e^{\frac{1}{2^{1/4}}}$}
\opt{B}{$f\left(\dfrac\pi2\right)=e$}
\opt{C}{$f\left(\dfrac{9\pi}{4}\right)=-2e^{\frac{1}{2^{1/4}}}$}
\opt{D}{$f(2\pi)=-2$}
\end{qoptions}
\qn{5} Let $y=f(x)$ be a curve passing through (1, 1) such that triangle formed by coordinate axis and tangent at any point of the curve lies in the first quadrant has area 2 square unit then
\begin{qoptions}
\opt{A}{order and degree of differential equation are 1 and 2 respectively}
\opt{B}{A linear differential equation can be formed}
\opt{C}{one of the curve is hyperbola}
\opt{D}{Curve passes through point (0, 2)}
\end{qoptions}
\qn{6} Let $\mathrm A=\displaystyle\int_1^{e^2}\dfrac{\ln x}{\sqrt x}\,dx$, which of following statements is/are correct.
\begin{qoptions}
\opt{A}{$\mathrm A<2\left(e-\dfrac1e\right)$}
\opt{B}{$\mathrm A<(e-1)\left(2+\dfrac{1}{\sqrt e}\right)$}
\opt{C}{$\mathrm A=\displaystyle\int_0^2 t\sqrt{e^t}\,dt$}
\opt{D}{$2<\mathrm A<5$}
\end{qoptions}
\sect{Section-II}
\qn{1} If $\displaystyle\int\dfrac{x^4+1}{x\left(x^2+1\right)^2}dx=a\ell n|x|+\dfrac{b}{1+x^2}+C$, then $\displaystyle\lim_{x\to0}\left(\dfrac{a\sin10x+b\cos x-1}{x}\right)$ is equal to

(Where C is integration constant, a, b $\in$ R)
\qn{2} If the curve $y=f(x)$ satisfying $\displaystyle\int_1^x t^2f(t)\,dt=\dfrac{x^2}{2\sqrt3}+\int_x^1 f(t)\,dt$ is intersected by the line $y=k$ at atleast two points, then the complete set of values of k is given by $k\in(-a,a)-\{b\}$. Then the value of $\left(a\sqrt3+b\right)$ is equal to
\qn{3} The triangle $\triangle$ABC with altitude CD and point E on side AC has $\triangle$ADE $\cong\triangle$CBD.
\qfig{p1_math_s2_q03_1}
$\dfrac{\text{Area of }\triangle\mathrm{ABC}}{\text{Area of }\triangle\mathrm{ADE}}=\dfrac{a+\sqrt b}{c}$

where a, b, c are natural numbers such that H.C.F. (a, b, c) = 1. Then the value of $\dfrac{a+b}{c}$ is equal to
\qn{4} Let $f(x)=3+\displaystyle\int_0^x\dfrac{1+\sin t}{t^2+2}\,dt\ \ \forall\ x\in\mathbb{R}$ and P(x) is a second degree polynomial function of x such that P(0) = f(0), P$'$(0) = f $'$(0) and P$''$(0) = f $''$(0). Let `A' be the area enclosed by y = P(x), the y-axis and the tangent at ($-$2, P($-$2)). Then the value of `2A'.
\qn{5} Let `n' be a positive integer such that $\displaystyle\int_0^1 e^{2x}(x-1)^n\,dx=\dfrac{e^2-5}{4}$. Let $\mathrm I_1=\displaystyle\int_{\sec^2z}^{n-\tan^2z}x\,f\left(x((n+1)-x)\right)dx$ and $\mathrm I_2=\displaystyle\int_{\sec^2z}^{n-\tan^2z}f\left(x((n+1)-x)\right)dx$ where `f' is a continuous function and $\mathrm Z\in\mathrm R-\left\{(2k+1)\dfrac\pi2\right\}$, where $k\in\mathrm R$ then the value of $\dfrac{\mathrm I_1}{\mathrm I_2}$ is equal to
\qn{6} The solution of the differential equation $e^{y^2}(y\,dy+x\,dx)=\left(x\sin x^2\right)dx$ with the condition y = 0 at x = 0 is $e^{y^2}=\dfrac{\sin\left(x^2\right)-\cos\left(x^2\right)}{2}+\dfrac32f(x)$. Then the value of $f\left(\sqrt{\ln2}\right)$ is equal to
\sect{Section-III}
\qn{1} If $f(x)$ is a continuous and differentiable function such that $\displaystyle\sum_{r=0}^{10}\left(f(x+r)-\left|e^x-r-1\right|\right)=0$ and $\displaystyle\int_0^{11}f(x)\,dx=a+b\ln2-ce$ where a, b, c $\in$ N . Then $\left|a-b^3-c\right|$ is
\qn{2} If a curve satisfying differential equation $\left(y^3-2x^2y\right)dx+\left(2xy^2-x^3\right)dy=0$ and passing through (1, 2). Then $[|y(2)|]$ is, where [ ] = Greatest integer function.
\qn{3} Let $y=f(x)\ \forall\ x\in\mathbb{R}$ satisfying $x.\left(y'+y''\right)=-y'-e^{-x}$ and $y(0)=-1$, then number of solution of $|f(x)|=|x|$ is
\qn{4} If area enclosed by curve $f(x)=\sin^{-1}\left(\dfrac{2|x|}{1+x^2}\right)+\sin^{-1}\left(\dfrac{|x|}{\sqrt{1+x^2}}\right)$ and $g(x)=\left|\tan^{-1}x\right|$ from x = $-1$ to x = 1 is $\pi-\ln b$. Then value of $[b-\pi]$ is, where [ ] is greatest integer function.
\qn{5} Let $\mathrm I_n=\displaystyle\int_0^{\pi/4}\tan^nx\,dx$ (n = 0, 1, 2, 3, 4, .....) and $\mathrm S_n=\displaystyle\sum_{n=0}^{n}\left(\mathrm I_n\mathrm I_{n+1}+\mathrm I_n\mathrm I_{n+3}+\mathrm I_{n+1}\mathrm I_{n+2}+\mathrm I_{n+2}\mathrm I_{n+3}\right)$. The value of $\displaystyle\lim_{n\to\infty}\mathrm S_n$ is
\qn{6} If $\displaystyle\int\dfrac{x(x+1)\left(2x^2-x+1\right)}{\left(x^3+x^2+x-1\right)^3}dx=\dfrac{1}{\mathrm Af^2(x)}+c$ with f(1) = 2. then $[f(2)]$ is (where [ ] is greatest integer function and A is real constant and c is constant of integration)
\end{multicols}
\end{document}
